% !TEX root = ../manuscript.tex

\section{Discussion}

\subsection{Implications for Site Selection, Operations, and Monitoring}

\todo{Translate the ensemble results into implications for screening faulted
storage sites, selecting injection rates and locations, and prioritizing
monitoring of the fault, top seal, and upward-migration pathways. Distinguish
robust implications from field-specific observations.}

\subsection{Value and Interpretation of Hierarchical Reduction}

The full uncertainty space combines 162 geologies, six window-specific
empirical libraries, multiple cross-window assignments, and along-strike
variability. Exhaustive reservoir simulation is therefore infeasible. The
hierarchical reduction retains actual joint permeability realizations,
dominant multimodal structure, and explicit coupling assumptions rather than
replacing the empirical distributions with component-wise averages or one
parametric model.

Window similarity groups identify distributions that can reasonably share a
state assignment in candidate coupled cases. They do not establish that
permeability realizations are physically correlated across windows. This
distinction is necessary because PREDICT is run independently by window and
does not identify the true cross-window joint distribution.

\subsection{Representation of Multiphase Fault Properties}

\todo{Discuss why spatially variable capillary pressure cannot generally be
replaced by a simple average of sand and clay reference curves, when entry-
pressure-branch medoids provide an acceptable simplification, and why a
representative relative-permeability curve is more defensible after endpoint-
consistent scaling. Relate these choices to reservoir-solver robustness.}

\subsection{Limitations and Transferability}

\todo{Address PREDICT structural assumptions, validation of clay-smear
placement, state and similarity thresholds, absence of observed cross-window
dependence, finite sampling and reservoir-simulation budgets, model and fluid-
property uncertainty outside the fault, and transferability beyond the
offshore Texas field.}
